%% Original source presentations supplement, not replace, the general CDCL proof.
\chapter{The Original GRASP and MiniSat Presentations}
\label{sec:p12-grasp-and-minisat-originals}

GRASP's 1999 journal article and the MiniSat contribution to SAT~2003,
published in 2004, have different native control structures and interfaces.
This chapter treats the publication-form sources
\cite{MarquesSilvaSakallah1999,EenSorensson2003}.
The finite-CNF first-UIP and watched-propagation construction in
Chapter~\ref{sec:p12-cdcl-uip-and-watched-literals} remains a separate proved
interface.  In particular, it does not by its name prove GRASP's same-frame
failure-driven assertions or MiniSat's arbitrary-constraint callbacks and
incremental-assumption return protocol.

\section{GRASP's antecedent assignments and cause sets}

\begin{hbdefinition}{def:p12-grasp-native-state}{The original clause and causality notation}
GRASP writes \(\varphi\) for the clause database, \(A\) for a partial
assignment, \(\nu(x)\) for the assigned value of a variable, and
\(\delta(x)\) for its level.  Its \(X=\{0,1\}\) is the unassigned
value marker, not the finite variable set denoted by \(X\) in the later
handbook presentation.  Preserve the source's literal convention
\(x^0=x\), \(x^1=\neg x\): \(x^{\nu(x)}\) is the literal made
\emph{false} by the assignment.  The plus sign in a displayed source clause
means disjunction, not integer addition.

For a unit reason clause \(\omega\) implying an assignment to \(x\),
\(A^\omega(x)\) is the set of assignments falsifying its other literals.
For example, from \(\omega=x+y+\neg z\),
\(A^\omega(x)=\{y=0,z=1\}\).  A decision has no antecedent assignments.
The implication graph \(I\) has assignment vertices, edges from the
antecedents to their implied assignment, and a conflict vertex \(\kappa\)
whose antecedents falsify its clause.  All antecedents precede their
consequence, making the graph acyclic
\cite[Secs.~2.1, 2.4; Fig.~1]{MarquesSilvaSakallah1999}.

For an analysis at level \(d\), partition the antecedents of a current-level
vertex into \(\Lambda(x)\), the assignments from levels below \(d\), and
\(\Sigma(x)\), those from level \(d\).  Treat the conflict vertex's level
as the maximum level of its antecedents.  Root facts are handled at level zero;
the maximum of an empty unit-reason set is taken as zero.  The latter is an
explicit boundary completion of the displayed maximum in source Eq.~(1),
not a claim that a maximum of an empty set exists without a convention.
\index{GRASP!antecedent assignments}\index{GRASP!cause sets}
\end{hbdefinition}

\begin{hbmethod}{def:p12-grasp-causes}{The standard backward cause construction}
Fix the current level \(d\).  Starting from the conflict, expand a
current-level implied vertex into its antecedents; retain lower-level
assignments without expanding them; retain a current-level decision as a
leaf.  In the source notation, the recursive operation is
\[
\operatorname{causes\_of}(x)=
\begin{cases}
 \{(x,\nu(x))\},&A(x)=\varnothing,\\
 \Lambda(x)\cup
 \displaystyle\bigcup_{(y,\nu(y))\in\Sigma(x)}
        \operatorname{causes\_of}(y),&\text{otherwise}.
\end{cases}
\]
Here \(A(x)\) suppresses the selected reason superscript; for the conflict
use its conflict antecedents rather than treating it as a Boolean variable.
An already empty conflict clause has an empty cause set and is an immediate
root refutation.  Put
\[
 A^{\omega_C}(\kappa)=\operatorname{causes\_of}(\kappa),\qquad
 \omega_C(\kappa)=
   \sum_{(x,\nu(x))\in A^{\omega_C}(\kappa)}x^{\nu(x)}.
\]
The sum is the disjunction of the assignments' opposite literals.  These are
the standard construction of source Eqs.~(2)--(4); the source separates it
from the UIP variants
\cite[Sec.~3; Eqs.~(2)--(4)]{MarquesSilvaSakallah1999}.
The recursion is over the actual antecedent graph, not a set of unrelated
variables that merely happens to have the same cardinality.
\end{hbmethod}

\begin{hblemma}{thm:p12-grasp-causes-sound}{proved=annex:p12-grasp-causes-sound}{Termination and semantic sufficiency of GRASP causes}
On every finite acyclic reason graph with sound selected reasons, the standard
cause construction terminates.  Every model satisfying its cause assignment
would reproduce the conflicting consequences; hence no model of the input
extends that assignment, and \(\omega_C(\kappa)\) is an input implicate.
The result includes conflicts following earlier failure-driven assertions:
their retained learned reasons must themselves be input implicates.  It does
not presume that such assertions are fresh decisions.
\end{hblemma}

\section{The native recursive search, assertions, and backtracking}

\begin{hbmethod}{def:p12-grasp-standard-engines}{Decide, Deduce, Diagnose, and Erase}
The original framework uses four engines
\cite[Sec.~2.5; Figs.~2--3, 5]{MarquesSilvaSakallah1999}.
\emph{Decide} selects an unassigned variable and value, or recognizes a
complete satisfying assignment.  The experimental decision heuristic counts
the clauses directly satisfied by each proposed variable/value assignment;
it is not Chaff's later activity heuristic.  \emph{Deduce} checks falsified
clauses before applying unit clauses, records each selected antecedent, and
continues until conflict or unit saturation.  \emph{Diagnose} constructs and
records \(\omega_C\), computes
\[
 \beta=\max\{\delta(x):(x,\nu(x))\in A^{\omega_C}(\kappa)\},
\]
and returns a conflict to an earlier frame if \(\beta\ne d\), recording
the new conflict's antecedents.  If \(\beta=d\), it leaves the learned
clause available for an assertion after \emph{Erase} removes the current
attempt.  An empty cause set is handled as a root conflict, not an undefined
maximum.  \emph{Erase} removes the assignments of the abandoned attempt,
not the fixed initial query assumptions.

Before the first decision, perform root Deduce on the fixed initial query,
reject a root conflict, and retain its fixed facts.  The recursive search
then first calls Decide.  In each attempt it calls Deduce.
If propagation succeeds, it searches a deeper frame.  A deeper conflict whose
\(\beta\) is below this frame is passed upward immediately.  A local
conflict, or one returned to this frame, is diagnosed.  If diagnosis points
below the frame, erase and pass it upward.  Otherwise erase this attempt and
repeat Deduce with the recorded clause, which forces the opposite branch.
The standard driver thus keeps the current frame for a failure-driven
assertion; it does not invariably jump to the second-highest learned-clause
level as the later first-UIP procedure does.

Level bookkeeping needs an explicit distinction in this source.  Eq.~(1)
gives an ordinary implied vertex the maximum antecedent level; the actual
Deduce pseudocode in Fig.~3 sets \(\delta(x)=d\), and Appendix Definition~4
classifies an FDA by antecedents strictly below its assigned level.  The
operational Search/Deduce presentation here uses the figure's current-frame
annotation, including FDA vertices.  The reason-max convention remains a
separate formula, not substituted during the same-frame FDA argument.
Likewise Fig.~2 starts its frame parameter at zero whereas §2.2 describes the
first decision as level one.  For the relative-query theorem below, root
facts and fixed assumptions occupy zero and the first decision frame is one;
this explicit initialization convention does not purport to verify an
unexplained zero-based executable by copying its wrapper token.
\end{hbmethod}

\begin{hbdefinition}{def:p12-grasp-fda-cdb}{FDA and highest-cause-level CDB}
If the cause set includes the present decision \(x=v@d\), its opposite
literal occurs in \(\omega_C\).  Removing that decision and its implication
sequence while retaining the earlier causes makes \(\omega_C\) unit and
forces \(x=1-v\).  This is a \emph{failure-driven assertion} (FDA), with
the learned clause as reason.  It is not an independent elective second
decision.  In the original current-frame convention its antecedents have
levels strictly below \(d\).

If a subsequent conflict is supported entirely by assignments below the
current frame, set \(\beta\) to their \emph{highest} level.  Every extension
of the prefix through \(\beta\) still falsifies the recorded conflict clause,
so deeper frames can be skipped.  At frame \(\beta\), the source diagnoses
that conflict again to produce an FDA or a still-earlier target.  This is
GRASP's standard conflict-directed backtracking (CDB), not an assertion that
the clause is already unit at \(\beta\)
\cite[Secs.~3.1.1--3.1.2; Eq.~(7)]{MarquesSilvaSakallah1999}.
\end{hbdefinition}

\begin{hbtheorem}{thm:p12-grasp-standard-search}{proved=annex:p12-grasp-standard-search}{The standard GRASP search with a declared root convention}
For finite CNF, sound complete unit propagation, standard cause diagnosis,
retained learned clauses, no restarts, and the fixed-root/current-frame
convention stated above, the recursive search terminates.  With fixed initial
assignment \(A_0\), success gives a total model extending \(A_0\), and
failure occurs exactly when there is no such model.  FDA and CDB skip only
assignments incompatible with an input implicate.  The theorem covers the
standard engine, not arbitrary engine substitutions, the red-clause variant,
or a compiled implementation of the published pseudocode.
\end{hbtheorem}

\begin{hbexample}{def:p12-grasp-native-conflict}{The complete native FDA/CDB conflict construction}
The original example uses
\[
\begin{array}{lll}
 \omega_1=\neg x_1+x_2,&\omega_2=\neg x_1+x_3+x_9,&
 \omega_3=\neg x_2+\neg x_3+x_4,\\
 \omega_4=\neg x_4+x_5+x_{10},&
 \omega_5=\neg x_4+x_6+x_{11},&\omega_6=\neg x_5+\neg x_6,\\
 \omega_7=x_1+x_7+\neg x_{12},&\omega_8=x_1+x_8,&
 \omega_9=\neg x_7+\neg x_8+\neg x_{13}.
\end{array}
\]
Retain \(x_9=0@1\), \(x_{12}=x_{13}=1@2\), and
\(x_{10}=x_{11}=0@3\).  Decide \(x_1=1@6\).  Then \(\omega_1\)
forces \(x_2=1\), \(\omega_2\) forces \(x_3=1\), \(\omega_3\)
forces \(x_4=1\), and \(\omega_4,\omega_5\) force \(x_5=x_6=1\),
falsifying \(\omega_6\).  Expanding the current-level antecedents yields
\[
 A^{\omega_C}(\kappa)=\{x_1=1,x_9=0,x_{10}=0,x_{11}=0\},\qquad
 \omega_C(\kappa)=\neg x_1+x_9+x_{10}+x_{11}.
\]
After erasure of the level-6 attempt, this clause forces the FDA \(x_1=0@6\).
Now \(\omega_8\) forces \(x_8=1\), \(\omega_7\) forces \(x_7=1\),
and \(\omega_9\) conflicts.  The FDA reason is expanded as well, giving
\[
 \omega_C(\kappa')=x_9+x_{10}+x_{11}+\neg x_{12}+\neg x_{13},\qquad
 \beta=3.
\]
There is no remaining cause from level 6.  The whole prefix through level 3
is excluded, so frames 6, 5 and 4 can be abandoned.  At level 3, further
diagnosis supplies the source's next FDA/backtrack decision
\cite[Figs.~1, 4; Eqs.~(5)--(7)]{MarquesSilvaSakallah1999}.
This is not the different \(x_{21}/x_{31}\) example of the later handbook manuscript.
\end{hbexample}

\section{The original UIP links and complete-product boundary}

\begin{hbmethod}{def:p12-grasp-uip-links}{UIP-local reasons with explicit endpoint order}
Let \(u_1,\ldots,u_k\) be the current-level UIPs in forward order from
the decision toward the conflict.  To explain a downstream vertex \(z\)
relative to upstream UIP \(u_i\), stop backward expansion at \(u_i\):
\(\operatorname{causes\_of}(u_i,u_i)=\{(u_i,\nu(u_i))\}\), and otherwise
use the same \(\Lambda/\Sigma\) recursion.  To obtain a link from \(u_i\)
to \(u_{i+1}\), negate the supporting assignments of
\(\operatorname{causes\_of}(u_{i+1},u_i)\) and append the literal
\(u_{i+1}^{1-\nu(u_{i+1})}\), which asserts the downstream value.  To
exclude the final conflict, negate
\(\operatorname{causes\_of}(\kappa,u_k)\).  Thus the source example yields
both \(\neg x_1+x_9+x_4\) and \(\neg x_4+x_{10}+x_{11}\).

This endpoint convention repairs the printed Eq.~(9): with Eq.~(8)'s argument
order, its displayed cause arguments and appended endpoint do not give the
stated links.  Merely retaining the stopping UIP and appending its own true
literal yields a tautology, not the desired implication.  The construction
above has its full semantic proof below; the original Eq.~(9) bytes remain
unchanged in the primary capture
\cite[Sec.~3.2.2; Eqs.~(8)--(9)]{MarquesSilvaSakallah1999}.
\end{hbmethod}

\begin{hblemma}{thm:p12-grasp-uip-links}{proved=annex:p12-grasp-uip-links}{Soundness of the corrected native UIP links}
Every link produced by the stated endpoint-ordered construction is an
implicate of the original input, and the terminal UIP conflict clause excludes
exactly a sufficient cause assignment.  No claim that every such link has a
universally smaller literal count than every other learned clause is needed
for soundness.  Choosing the first UIP alone is a different stopping policy.
\end{hblemma}

\begin{hbmethod}{def:p12-grasp-retention-policy}{The source's green/red width policy}
The source also distinguishes permanent green clauses of size at most a
fixed \(k\) from larger red clauses, which are removed once they have more
than one free literal.  This source-specific bounded diagnosis changes the
retention policy and is not covered by the standard-retained-search theorem
\cite[Sec.~3.2.1]{MarquesSilvaSakallah1999}.
\end{hbmethod}

\begin{hbdefinition}{def:p12-grasp-complete-product}{Prime implicates and the complete product}
A non-tautological clause is a prime implicate of finite CNF \(\varphi\)
when it is entailed by \(\varphi\) and no proper literal subset is entailed.
Let \(P(\varphi)\) be the set of all such clauses.  If \(\varphi\) is
unsatisfiable, the empty clause is the unique prime implicate under this
definition; if \(\varphi\) is identically true, \(P(\varphi)\) is empty.
This explicitly corrects the source appendix's phrase that the false function
has no prime implicates while simultaneously placing the empty clause in its
complete product.  Also, a prime subclause \(\omega\subseteq\psi\)
implies \(\psi\), not the reverse implication printed in the satisfiable
case of its proof
\cite[Appendix, Thm.~3, pp.~519--520]{MarquesSilvaSakallah1999}.
\end{hbdefinition}

\begin{hbtheorem}{thm:p12-grasp-complete-product}{proved=annex:p12-grasp-complete-product}{Conflict-free decisions with all prime implicates}
For satisfiable finite \(\varphi\), every consistent partial assignment
that is unit-saturated and nonconflicting on \(P(\varphi)\) extends to a
model.  Choosing either value of a still-unassigned variable and propagating
cannot produce a conflict.  For an unsatisfiable input, the empty prime
implicate gives an immediate conflict.  The satisfiable claim is for a
compatible initial prefix, not for externally imposed contradictory query
assumptions, and implies no small-size bound for \(P(\varphi)\).
\end{hbtheorem}

\section{MiniSat's actual 2004 callback and result contracts}

\begin{hbdefinition}{def:p12-minisat-original-interface}{The published types, state, and external API}
The original paper uses \texttt{var}, \texttt{lit}, \texttt{Vec},
\texttt{Queue}, and \texttt{lbool}; \(\bot\) denotes \emph{undefined},
not Boolean false.  The lifted values are
\(\mathrm{True}_{\bot},\mathrm{False}_{\bot},\bot\), with distinct
undefined literal and variable sentinels.  The public interface creates
variables with \texttt{newVar}, adds clauses or declared Boolean constraints,
calls \texttt{simplifyDB}, and calls \texttt{solve(assumptions)}; a successful
solve exposes a \texttt{model} vector
\cite[Secs.~2, 4.1; Figs.~1--3]{EenSorensson2003}.

The state fields are the problem-constraint vector \texttt{constrs}, learned
clause vector \texttt{learnts}, variable/clause scores and bump/decay factors,
\texttt{order}, per-true-literal \texttt{watches}, per-variable
\texttt{undos}, \texttt{propQ}, \texttt{assigns}, \texttt{trail}, level-boundary
indices \texttt{trail\_lim}, \texttt{reason}, \texttt{level}, and
\texttt{root\_level}.  The last field separates temporary query assumptions
from search decisions.  These names and types are native implementation
objects, not new inference rules.

A false return from addition or top-level simplification leaves the source's
solver object unusable: it has no stored Boolean ``already inconsistent''
flag.  Before nonempty-assumption solve, top-level \texttt{simplifyDB} must
therefore establish a usable base.  A false \texttt{solve} return then means
UNSAT \emph{under those assumptions}; cancellation restores a usable
assumption-free state.  It does not assert that the permanent base is UNSAT
\cite[Sec.~2, pp.~503--504; Fig.~8]{EenSorensson2003}.
\end{hbdefinition}

\begin{hbmethod}{def:p12-minisat-original-callbacks}{The native constraint protocol}
Each Boolean constraint implements \texttt{remove}, \texttt{propagate}, and
\texttt{calcReason}; \texttt{simplify} and \texttt{undo} have defaults.
Construction at top level registers all relevant watch events and enqueues
the unit information available there.  Removal detaches registrations before
disposing the object.  Propagation is invoked by a newly true watched literal;
it reregisters itself in the same or another watch list, enqueues justified
consequences, and reports a genuine conflict with a false return.

\texttt{calcReason(p)} returns true assignment literals supporting the
propagated \(p\); \texttt{calcReason(}\(\bot_{\mathrm{lit}}\)\texttt{)} returns assignments
supporting a conflict.  The source guarantees the callback is evaluated in
the relevant earlier assignment state.  A nonclausal constraint that changes
internal state during propagation may register an undo callback; reverse-trail
cancellation invokes it with that earlier state restored.  The no-undo
property for clauses is not automatically a no-undo property for every
Boolean constraint type \cite[Sec.~4.2]{EenSorensson2003}.

For the operational return convention, \texttt{simplify} returns true only
when the constraint is satisfied and may be deleted, and false when it remains
after in-place simplification.  This follows the actual Clause code in Fig.~4
and its caller in Fig.~9.  The prose on p.~508 reverses those two meanings;
the present protocol corrects that discrepancy rather than declaring both
variants simultaneously executable.  A stronger simplifier detecting a
conflict requires its result to be checked: the source search ignores that
result only for its particular BCP-strength simplifier
\cite[pp.~508, 510, 515, 517]{EenSorensson2003}.
\end{hbmethod}

\begin{hbdefinition}{def:p12-minisat-original-reasons}{Semantic obligations of returned reason sets}
For a Boolean constraint \(K\), a propagated literal \(p\), and a finite
set \(R\) of previously true assignment literals, a sound callback requires
\[
 K\models\bigl((\bigwedge R)\to p\bigr).
\]
For a conflict reason it requires \(K\models\neg\bigwedge R\).
Thus the corresponding clause is the disjunction of the complements of
\(R\), with \(p\) appended in the propagation case.  This semantic
obligation is required even when \(K\) is not itself a clause.  Registered
events must also expose violated constraints at a total assignment; a
callback that is merely sound when called but is never called on the needed
event does not justify SAT.  These are the proof obligations of the source
extension interface, not an implementation supplied by its type name.
\end{hbdefinition}

\begin{hbtheorem}{thm:p12-minisat-original-reason-bridge}{proved=annex:p12-minisat-original-reason-bridge}{The nonclausal reason-to-clause bridge}
Under these reason, chronological-support, event-coverage and undo contracts,
every clause learned by the native callback analysis is entailed by the
permanent Boolean constraints.  A total conflict-free return satisfies them;
a root-query conflict excludes all models satisfying the supplied assumptions.
This proves the semantic bridge to the general CDCL resolution argument,
not correctness or termination of an arbitrary new callback implementation.
\end{hbtheorem}

\section{Published propagation, analysis, and query execution}

The original \texttt{Clause} stores \texttt{learnt}, \texttt{activity}, and
\texttt{lits}; the two watches occupy \texttt{lits[0]} and \texttt{lits[1]}.
Its propagation swaps the newly false watched literal into position 1,
accepts immediately if position 0 is true, otherwise searches positions 2
onward for a nonfalse replacement.  On failure it reregisters the old watch
and enqueues position 0, which either propagates or conflicts.  On a conflict,
the outer list traversal restores all unvisited list references.  Original
\texttt{enqueue} immediately writes the assignment, reason and current level;
an opposite enqueue clears \texttt{propQ} and returns false, while an already
true enqueue does nothing \cite[Figs.~4, 6]{EenSorensson2003}.
The abstract watched theorem can be applied only with these exact event and
attachment conditions, as in Theorem~\ref{thm:p12-cdcl-uip-watch-correct}.

The original constructor normalizes input clauses at top level, detects a
true literal or tautology, removes false/duplicate occurrences, rejects an
empty residual, and enqueues a unit residual.  A learned nonunit clause has
its asserting literal first and a highest-level nonasserting watch second.
The published Fig.~5 transfers \texttt{ps} into \texttt{c.lits} and then
still indexes \texttt{ps}; the abstract vector transfer semantics are not
specified by the printed API.  A concrete implementation must therefore
specify the ownership boundary.  With a source-clearing transfer, subsequent
watch selection and bumping must use \texttt{c.lits}, not the emptied vector.
This is an explicit implementation dependency, not a claim that the literal
pseudocode has been compiled and verified.

\begin{hbexample}{def:p12-minisat-original-reason-example}{The published callback-reason expansion}
The publication's own learning example makes the reason types concrete.
Here \(u,v,x,y,z\) are Boolean variables, not frame indices or assigned-value
parameters of the earlier GRASP presentation.
If \(\{x,y,z\}\) is false, its conflict reason consists of the true
assignments \(\neg x,\neg y,\neg z\).  If a constraint explains
\(\neg x\) by true assignments \(u,\neg v\), it entails
\(\neg u+v+\neg x\).  Expanding that callback reason replaces
\(\neg x\) in the conflict reason by \(u,\neg v\), producing the
learned clause \(\neg u+v+y+z\).  A reason set is therefore a conjunction
of true assignments, not the disjunction returned as the learned clause
\cite[Sec.~4.4, p.~513]{EenSorensson2003}.
\end{hbexample}

The native analysis uses \texttt{seen}, \texttt{counter}, a reason vector and
a backward trail scan.  Initially \(p=\bot_{lit}\) requests the conflict
reason.  Newly seen current-level variables increase the counter; positive
older levels contribute complemented literals to the result and determine
the backjump maximum.  Root-zero literals are omitted only because they are
already permanent consequences.  The scan undoes trail entries until it
reaches a seen literal, takes its stored reason, decrements the counter, and
repeats until one UIP remains.  Its complement becomes the first learned
literal.  The caller completes cancellation to the larger of the computed
backjump and \texttt{root\_level}, then records and enqueues the assertion
\cite[Fig.~7]{EenSorensson2003}.

A learned unit is not allocated as a Clause object or inserted into
\texttt{learnts}; the constructor enqueues it as a fact.  Under temporary
assumptions its computed backjump may be zero but cancellation is clamped to
the query root, so that fact is enqueued there.  Cancelling the assumptions
after solve can erase it, and a later query may have to derive it again.
The original explicitly describes this unit-storage boundary; it does not
implement permanent storage of every learned unit
\cite[p.~514]{EenSorensson2003}.  Forgetting an entailed
fact cache is not a claim that the permanent input changed its models.

\texttt{search} first propagates.  A conflict at the query root returns
\(\mathrm{False}_{\bot}\); other conflicts invoke analysis, cancellation,
recording and activity decay.  Without a conflict it may simplify at absolute
level zero and reduce unlocked learned clauses.  A total assignment is copied
to \texttt{model}; exceeding the conflict budget cancels to the query root
and returns \(\bot\), not UNSAT.  The outer \texttt{solve} installs each
assumption with propagation, fixes \texttt{root\_level}, repeats bounded
searches, and eventually cancels temporary assumptions.  Its printed strategy
starts with conflict budget 100 and learned limit \(\mathrm{nConstraints}/3\),
then multiplies them by 1.5 and 1.1; variable/clause decay parameters are
0.95 and 0.999.  These are the original selected strategy constants, not
general complexity bounds \cite[Fig.~8]{EenSorensson2003}.

\begin{hbtheorem}{thm:p12-minisat-original-query}{proved=annex:p12-minisat-original-query}{Root simplification and assumption-relative conclusions}
Let \(F\) be the permanent Boolean constraint conjunction and \(G\) a
finite conjunction of root facts entailed by \(F\).  If each retained
constraint is simplified equivalently under \(G\), and only constraints
already true under \(G\) are removed, the represented conjunction
\(G\land F'\) has exactly the models of \(F\).  With temporary assumption
conjunction \(H\), a checked model represents a model of \(F\land H\),
and a sound query-root conflict proves \(F\land H\) unsatisfiable.
Neither proves \(F\) unsatisfiable merely because \(H\) was inconsistent
with it.  A budget return of undefined carries neither conclusion.
\end{hbtheorem}

\begin{hbexample}{def:p12-minisat-original-assumptions}{A native return-value boundary}
The permanent clause \(x+y\) is SAT, for example at \(x=1,y=0\).
A solve call with assumptions \(\neg x,\neg y\) conflicts and returns
false \emph{under assumptions}; it does not refute the permanent clause.
After cancellation a solve call with assumption \(x\) still has a model.
If a caller instead ignores a false addition of the contradictory units
\(x\), \(\neg x\), the source declares that object unusable.  Calling
assumption solve on it and interpreting the next false as merely temporary
violates the source's precondition; no theorem here validates that misuse.
\end{hbexample}

The 17-page published contribution and the 23-page extended version~1.2 are
not title aliases.  The original first-page footnote points to an extended
version; its nine grouped code figures have different numbering from the
extended presentation, which separately includes extra ADT exposition,
exercises and a Satzoo appendix.  The original false-enqueue queue-clearing
step differs from extended~1.2 Fig.~9, whose enqueue false branch lacks that
local clear while its outer propagation clears the queue on callback failure.
The original Simplify prose reverses the flags, while the extended prose
uses the flags required by its caller.  These verified local differences
must not be erased by substituting later text.  The shared
first-UIP/watch case structure explains a bounded relationship to the general
proof chapter; it is not a proof that all program states, API flags, wrapper
conventions or empirical results of these versions are identical.
